\section{Introduction}
\label{sec:intro}

\IEEEPARstart{S}{leep} loss impairs performance in ways that matter for safety. In a classic comparison, 17~h of sustained wakefulness reduced performance on a tracking task to the level seen at a blood alcohol concentration of 0.05\%, and 24~h to the level seen at 0.10\%~\cite{dawson1997fatigue}. People are also poor judges of the impairment that builds up over days of short sleep: under chronic sleep restriction, lapses on the psychomotor vigilance test (PVT) keep increasing while sleepiness ratings level off~\cite{vandongen2003cumulative}, and individual differences in subjective and objective vulnerability to sleep loss are stable but unrelated~\cite{leproult2003individual}.

Three kinds of tool address parts of this problem. \emph{Biomathematical fatigue models} predict sleepiness or performance from the sleep--wake history and the time of day~\cite{borbely1982two,akerstedt1997three,akerstedt2004predictions,hursh2004fatigue,mallis2004summary,dawson2011modelling}. They need no sensor, and some of them are individualised from repeated performance tests~\cite{olofsen2004nonlinear,vandongen2007optimization,rajaraman2008individualized,liu2017real,reifman2019twobalert}. \emph{Ocular measures} such as PERCLOS and blink duration track drowsiness~\cite{wierwille1994research,dinges1998evaluation,caffier2003experimental,schleicher2008blinks}, and blinks can be detected in ordinary video~\cite{soukupova2016real}, but most systems map them to drowsiness with one rule for everyone, although the link between eyelid signs and sleepiness differs between people~\cite{ingre2006subjective}. \emph{Validated questionnaires}, such as the Karolinska Sleepiness Scale (KSS)~\cite{akerstedt1990subjective}, capture how sleepy a person feels and track performance and EEG changes~\cite{kaida2006validation,akerstedt2014subjective}, but they can be asked only a few times a day.

This paper asks whether an ordinary laptop can combine all three: keep a continuously updated estimate of how sleepy a person feels, state its uncertainty, and project it forward, with video kept on the device. We call the result a personal digital twin in the sense of a person-specific model that is updated from that person's data and can be queried for forecasts and what-if scenarios~\cite{grieves2017digital,corralacero2020digital}. The model itself is a small state-space model, and we use that term in the technical sections.

The difficulty lies in combining sources that disagree. A webcam or reaction-time signal can read a particular person consistently high or low. If such a signal is averaged with the person's ratings, it pulls the estimate away from what the person reports. BlinkTwin therefore gives each passive signal a forward measurement model with its own personal \emph{anchor}: a slowly drifting bias state, in the signal's own units, that absorbs where this person's signal sits relative to the population. Check-ins set the level of the estimate, and passive signals contribute mainly through their changes. Bias states in a Kalman filter are a standard technique~\cite{friedland1969treatment}; the contribution here is to apply them to self-rated sleepiness and to test the result.

Because the KSS is the reference, BlinkTwin estimates \emph{subjective} sleepiness. Given the dissociation described above, the application also shows the latest reaction test read on the population scale, and it flags a possible gap only when the person's own reaction times slow over a day or a week while their ratings do not.

The short answer from public data and simulation is that check-ins and the biomathematical prior do most of the work. Passive signals add a little, and only when the model estimates each person's bias in each signal; a filter that merely allows for those biases gains nothing. The website's camera pipeline measures blinks on laboratory video but not PERCLOS, and the passive signals' benefit has yet to be shown in real use. We report these limits alongside the method because they decide what the planned study has to test.

\paragraph*{Contributions}
\begin{enumerate}
  \item An open, on-device implementation that joins the published Three-Process Model of alertness~\cite{akerstedt2004predictions,akerstedt2008accounting,ingre2014validating} with KSS check-ins and forward models of a 3-min PVT (PVT-B), blink duration and PERCLOS in one Kalman filter, with a personal anchor for each passive signal (Sections~\ref{sec:system} and~\ref{sec:model}).
  \item An evaluation on public data: a leave-one-subject-out replay of the model on DROZY~\cite{massoz2016ulg}; a test of the website's camera pipeline on DROZY's near-infrared video against the electro-oculogram (EOG), which exposed a failure of PERCLOS that we then guarded against; and tests of the offset premise and of the prior on OpenNeuro ds004902~\cite{xiang2024resting} (Section~\ref{sec:realdata}).
  \item A simulation study, reported following Morris et al.~\cite{morris2019using}, with a consider-covariance comparator~\cite{schmidt1966application}, a known-parameter reference, scenarios that break the model's assumptions, and a sensitivity analysis of the filter settings (Section~\ref{sec:simulation}).
  \item A planned protocol for a prospective validation study, to be registered before data collection (Section~\ref{sec:protocol}).
\end{enumerate}

BlinkTwin is a self-monitoring and research tool for desk work. It is not a medical device, it must not be used to decide whether someone is fit to drive or operate machinery, and the absence of an alert does not indicate alertness.
